% WARD's modular pipeline processes RGB frames through parallel SegFormer and Depth Anything V2 modules. 
% % These outputs are fused via our proposed continuous collision-corridor risk model to generate per-obstacle risk scores. To ensure temporal stability, these scores are smoothed by an extended SORT\cite{wojke2017simpleonlinerealtimetracking} tracker before being mapped to discrete warning levels.
% These outputs are fused via our continuous collision-corridor risk model to generate per-obstacle risk scores. To ensure stability, an extended SORT\cite{wojke2017simpleonlinerealtimetracking} tracker smooths these scores before mapping to discrete warning levels.

% \subsection{Segmentation and Calibrated Depth}
% We adopt SegFormer (MiT-B0) trained on the LaRS~\cite{zust2023larsdiversepanopticmaritime} dataset as our segmentation backbone. To enable instance separation from a semantic-only backbone, we extend it with a dual-head design: a 9-class semantic head and an auxiliary obstacle-boundary head, the latter used to split adjacent same-class instances during connected-component extraction. Depth is estimated using the Metric-Outdoor variant of Depth Anything V2, fine-tuned from the foundation model on synthetic Virtual KITTI data~\cite{gaidon2016virtualworldsproxymultiobject}. To correct scale bias from road-centric pre-training without retraining on scarce maritime data, we apply a single global scaling factor $s_{cal}=0.4$ calibrated once against radar ground truth from the WaterScenes dataset~\cite{yao2024waterscenesmultitask4dradarcamera} according to 
% $d_\mathrm{cal} = s_{cal} \cdot d_\mathrm{raw}$, 
% where $d_{\mathrm{raw}}$ is the raw output from the Depth Anything V2 model and $d_{\mathrm{cal}}$ represents the calibrated metric depth in meters. This adjustment restores near-metric accuracy in the critical $0$--$20$\,m close-up range, which is the primary zone for maritime collision warning. %For each segmented instance, a conservative depth is computed as the 5th percentile of its per-pixel depth distribution.

% \subsection{Risk Model and Extended Tracking}
% Based on the Responsibility-Sensitive Safety (RSS) framework~\cite{shalevshwartz2018formalmodelsafescalable}, we define a virtual forward collision corridor of width $W_\mathrm{boat} + 2M_\mathrm{lat}$ centered on the ASV heading, where $W_\mathrm{boat}$\ is the ASV beam width and $M_\mathrm{lat}$\ is the lateral safety margin. For an obstacle with forward distance $d_\mathrm{fwd}$ and lateral excess $e_\mathrm{lat}$ outside this corridor, we first compute a static risk baseline:
% \begin{equation}
%   R_\mathrm{base} = w_\mathrm{class} \cdot \exp\!\left(-d_\mathrm{fwd}/D_\mathrm{safe}\right) \cdot \exp\!\left(-e_\mathrm{lat}/L_\mathrm{safe}\right),
%   \label{eq:risk_base}
% \end{equation}
% where $w_\mathrm{class}$ is the semantic hazard weight, $D_\mathrm{safe}$\, is the forward decay constant, and $L_\mathrm{safe}$\, is the lateral decay constant. This baseline is amplified by a closing-velocity factor:
% \begin{equation}
%   R = R_\mathrm{base} \cdot \bigl(1 + \alpha_v \cdot \min(v_\mathrm{closing} / V_\mathrm{ref},\; 2)\bigr),
%   \label{eq:risk}
% \end{equation}
% where $v_\mathrm{closing}$ is the Kalman-filtered closing velocity, $V_\mathrm{ref}$\, is the reference speed, and $\alpha_v$ is the amplification factor. The continuous score $R$ is mapped to three discrete warning levels, SAFE ($R < 0.15$), CAUTION ($0.15 \le R < 0.50$), and IMMEDIATE($R \ge 0.50$), via threshold hysteresis with separate enter/exit values to suppress noise-induced flicker.
% To mitigate monocular depth noise, we extend the SORT tracker to a 9-D Kalman state $\mathbf{x} = [c_x, c_y, s, r, d, \dot{c}_x, \dot{c}_y, \dot{s}, \dot{d}]^\top$, where $(c_x, c_y)$, $s$, and $r$ denote the bounding-box center, scale (area), and aspect ratio; $d$ is the calibrated metric depth; and dotted variables represent their respective rates of change.
% The filter optimally balances depth observation noise against a constant-velocity process model. This yields smoothed estimates of the closing velocity, $v_{\mathrm{closing}} = -\dot{d} \cdot f_{\mathrm{fps}}$, and the Time-to-Collision, $TTC = d / v_{\mathrm{closing}}$, where $f_{\mathrm{fps}}$ is the system frame rate. These parameters feed directly into the risk model in Eq.~(\ref{eq:risk}) without requiring post-hoc exponential moving averages or windowing.

WARD's modular pipeline processes RGB frames through parallel SegFormer and Depth Anything V2 modules. These outputs are fused via our continuous collision-corridor risk model to generate per-obstacle risk scores. To ensure stability, an extended SORT tracker smooths these scores before mapping to discrete warning levels.

\subsection{Maritime Semantic Segmentation}
We adopt SegFormer (MiT-B0) trained on the LaRS~\cite{zust2023larsdiversepanopticmaritime} dataset as our segmentation backbone. To enable instance separation from a semantic-only backbone, we extend it with a dual-head design: a 9-class semantic head and an auxiliary obstacle-boundary head, the latter used to split adjacent same-class instances during connected-component extraction.

LaRS provides both semantic and instance-level annotations under conditions such as water reflections, sun glitter, and partially submerged obstacles. The MiT-B0 encoder was selected among the architectures we benchmarked for its accuracy--latency trade-off. The nine semantic classes, Static Obstacle, Water, Sky, Boat, Buoy, Swimmer, Animal, Float, and Other, are obtained by merging the LaRS panoptic categories. The two heads are trained jointly with a combined cross-entropy, Dice, and boundary binary cross-entropy loss, with class-specific weights (e.g., $5\times$ for Swimmer, $4\times$ for Buoy) to counter the severe class imbalance among hazard categories. At inference, the segmentation mask and calibrated depth map are fused: boundary pixels above a threshold $\tau_b=0.4$ are first suppressed to split adjacent same-class instances, and connected-component analysis is then applied to every non-background class (excluding Water and Sky) to isolate contiguous obstacle regions, discarding regions smaller than 300\,px as segmentation noise. Each surviving instance is assigned a representative depth equal to the 5th percentile of its pixel-wise depth distribution, a conservative estimate of its nearest point, and a normalized lateral offset $\ell \in [-1,1]$ relative to the image center.

\subsection{Monocular Depth Estimation}
Depth is estimated using the Metric-Outdoor variant of Depth Anything V2, fine-tuned from the foundation model on synthetic Virtual KITTI data~\cite{gaidon2016virtualworldsproxymultiobject}. To correct scale bias from road-centric pre-training without retraining on scarce maritime data, we apply a single global scaling factor $s_{cal}=0.4$ calibrated once against radar ground truth from the WaterScenes dataset~\cite{yao2024waterscenesmultitask4dradarcamera} according to 
$d_\mathrm{cal} = s_{cal} \cdot d_\mathrm{raw}$, 
where $d_{\mathrm{raw}}$ is the raw output from the Depth Anything V2 model and $d_{\mathrm{cal}}$ represents the calibrated metric depth in meters. This adjustment restores near-metric accuracy in the critical $0$--$20$\,m close-up range, which is the primary zone for maritime collision warning.

The scaling factor is fit using the sample subset of WaterScenes, which provides synchronized RGB and 4D-radar streams; radar points are projected onto image coordinates via the dataset's calibration to yield per-pixel metric ground truth. Even after calibration, accuracy degrades beyond 20\,m, and frame-to-frame estimates retain residual noise (median inter-frame variation of 0.28\,m on our evaluation sequences). Instead of pursuing tighter calibration, we design the downstream modules to tolerate this noise: the risk model relies on a conservative per-instance depth percentile rather than raw pixel depth, the tracker below smooths depth over time, and warning-level hysteresis absorbs the remaining flicker.

\subsection{Multi-Object Tracking}
To mitigate monocular depth noise, we extend the SORT tracker to a 9-D Kalman state $\mathbf{x} = [c_x, c_y, s, r, d, \dot{c}_x, \dot{c}_y, \dot{s}, \dot{d}]^\top$, where $(c_x, c_y)$, $s$, and $r$ denote the bounding-box center, scale (area), and aspect ratio; $d$ is the calibrated metric depth; and dotted variables represent their respective rates of change.
The filter optimally balances depth observation noise against a constant-velocity process model. This yields smoothed estimates of the closing velocity, $v_{\mathrm{closing}} = -\dot{d} \cdot f_{\mathrm{fps}}$, and the Time-to-Collision, $TTC = d / v_{\mathrm{closing}}$, where $f_{\mathrm{fps}}$ is the system frame rate. These parameters feed directly into the risk model in Eq.~(\ref{eq:risk}) without requiring post-hoc exponential moving averages or windowing.

Data association uses the Hungarian algorithm on an IoU cost matrix between predicted and observed boxes, with a minimum IoU of 0.20; unmatched detections spawn new tracks, tracks left unmatched for five consecutive frames are deleted, and a track is confirmed after two consecutive hits. Within that five-frame grace window, the filter continues to propagate depth and bounding-box motion forward under the constant-velocity model even without a matching detection, so a brief occlusion degrades the closing-velocity estimate gracefully rather than resetting it outright. The normalized lateral offset of each track is separately smoothed with an exponential moving average ($\alpha=0.4$), which keeps small shifts in the segmentation mask from producing abrupt changes in the lateral term of the risk model below.

\subsection{Collision-Risk Modeling}
Based on the Responsibility-Sensitive Safety (RSS) framework~\cite{shalevshwartz2018formalmodelsafescalable}, we define a virtual forward collision corridor of width $W_\mathrm{boat} + 2M_\mathrm{lat}$ centered on the ASV heading, where $W_\mathrm{boat}$\ is the ASV beam width and $M_\mathrm{lat}$\ is the lateral safety margin. For an obstacle with forward distance $d_\mathrm{fwd}$ and lateral excess $e_\mathrm{lat}$ outside this corridor, we first compute a static risk baseline:
\begin{equation}
  R_\mathrm{base} = w_\mathrm{class} \cdot \exp\!\left(-d_\mathrm{fwd}/D_\mathrm{safe}\right) \cdot \exp\!\left(-e_\mathrm{lat}/L_\mathrm{safe}\right),
  \label{eq:risk_base}
\end{equation}
where $w_\mathrm{class}$ is the semantic hazard weight, $D_\mathrm{safe}$\, is the forward decay constant, and $L_\mathrm{safe}$\, is the lateral decay constant. This baseline is amplified by a closing-velocity factor:
\begin{equation}
  R = R_\mathrm{base} \cdot \bigl(1 + \alpha_v \cdot \min(v_\mathrm{closing} / V_\mathrm{ref},\; 2)\bigr),
  \label{eq:risk}
\end{equation}
where $v_\mathrm{closing}$ is the Kalman-filtered closing velocity, $V_\mathrm{ref}$\, is the reference speed, and $\alpha_v$ is the amplification factor. The continuous score $R$ is mapped to three discrete warning levels, SAFE ($R < 0.15$), CAUTION ($0.15 \le R < 0.50$), and IMMEDIATE($R \ge 0.50$), via threshold hysteresis with separate enter/exit values to suppress noise-induced flicker.

Concretely, $W_\mathrm{boat}=2$\,m and $M_\mathrm{lat}=1$\,m, giving the lateral excess $e_\mathrm{lat} = \max(0,\ |d_\mathrm{lat}| - (W_\mathrm{boat}/2+M_\mathrm{lat}))$, where $d_\mathrm{lat} = \ell \cdot d_\mathrm{fwd} \cdot \tan(\theta_\mathrm{HFOV}/2)$. The decay constants are $D_\mathrm{safe}=12$\,m and $L_\mathrm{safe}=2$\,m; the per-class weight $w_\mathrm{class}$ runs from 1.0 for Boat down to 0.7 for Buoy and 0.3 for Static Obstacle; and the velocity term uses $V_\mathrm{ref}=3$\,m/s with $\alpha_v=1.0$. The multiplicative form requires both proximity and lateral alignment to be threatening at once, while the velocity term penalizes an approaching obstacle more heavily than a stationary or receding one at the same range. An obstacle directly ahead at $d_\mathrm{fwd}=0$ returns $R=w_\mathrm{class}$, and with the velocity factor included the score can reach $w_\mathrm{class} \cdot (1+2\alpha_v)$, so a high-weight obstacle closing rapidly can exceed 1.0 rather than saturate at a nominal ceiling. In practice, the enter/exit hysteresis thresholds for IMMEDIATE are set at 0.50 and 0.38, respectively, so a warning only drops once the score falls comfortably below the value that triggered it.
